
\DocumentMetadata{tagging=on,
	tagging-setup={math/setup=mathml-SE},
	pdfstandard=ua-2,
	lang=en-US
}
\documentclass{article}

%Math Libraries
\usepackage{multirow, longtable, amsmath, amssymb, amsthm}
\usepackage[mathscr]{euscript}

%Formatting
\usepackage[margin = 1 in]{geometry}
\usepackage{indentfirst}
\usepackage{graphicx}
\usepackage{fancyhdr}
\usepackage{tabularx}
\usepackage{cite}

%Diagram Packages
\usepackage{tikz}
\usetikzlibrary{automata,positioning,arrows}

%Standard Commands
\newcommand{\mm}{\mathscr}
\newcommand{\B}{{\mathcal{B}}}
\newcommand{\A}{{\mathcal{A}}}
\newcommand{\N}{{\mathbb{N}}}
\newcommand{\R}{{\mathbb{R}}}
\newcommand{\C}{{\mathbb{C}}}
\newcommand{\Q}{{\mathbb{Q}}}
\newcommand{\Z}{{\mathbb{Z}}}


\newcommand{\abs}[1]{{|{#1}|}}
\newcommand{ \norm }[1]{{ || {#1}||}}
\newcommand{\cl}[1]{{\overline{#1}}}
\newcommand{\pd}{\partial}
\newcommand{\ip}[2]{ \langle {#1},{#2}\rangle }
\newcommand{\ls}{\lesssim}
\newcommand{\gs}{\gtrsim}
\newcommand{\supp}{\text{supp}}

%Result Environment
\newcounter{result}[section]
\newenvironment{res}[1][]{\refstepcounter{result}\par\medskip
	\noindent \textbf{[\thesection.\theresult] #1} \rmfamily}{\medskip}

\newenvironment{defn}[1][]{\refstepcounter{result}\par\medskip \noindent \textbf{Definition [\thesection.\theresult] #1} \rmfamily}{\medskip}

%Proof Environment
\newenvironment{pf}{%
	\par%
	\medskip
	\leftskip=2em\rightskip=2em%
	\noindent\ignorespaces \textit{Proof:} \newline}{%
	\,\, \newline \textit{Q.E.D.}\par\medskip}

\title{Math 126 Summer 2026 HW 3}
\author{Lec 001, Ebbighausen}


\begin{document}
	\pagestyle{fancy}
	\fancyfoot{} 
	\fancyfoot[L]{Math 126}
	
	\maketitle
	
	\textbf{Problem 1:} Consider the example of a nonlinear diffusion equation 
	\begin{align*}
		\pd_{t} u - \Delta u^2 =0 \hspace{0.5cm} \{t \geq 0\} \times \R^{n} 
	\end{align*}
	
	Part A.) Assume that $u(t,x) = h(t)\phi(x)$. Find differential equations for $v$ and $\phi$ (we will solve in parts B. and C). 
	
	Part B.) Solve for $v(t)$. If we are given the initial condition $v(0) = a>0$, what happens to the solution at time $a$?
	
	Part C.) Assume $\phi$ is radial and a monomial in the radius, so $\phi(r) = r^{\alpha}$ for some $\alpha > 0$. What must $\alpha$ be to solve the equation from part A? 
	
	Part D.) Form the product solutions. 
	
	\textbf{Problem 2:} (Borthwick 5.7 Modified)
	
	The quantum energy levels of a harmonic oscillator in $\R^{n}$ are the eigenvalues of 
	\begin{align*}
		(-\Delta + |x|^{2}) \phi = \lambda \phi
	\end{align*}
	
	In the $n=1$ case, using an exponential term $\phi(x) = q(x)e^{-x^{2}/2}$ reduces to 
	\begin{align*}
		k qe^{-x^{2}/2} = -q''e^{-x^{2}/2} +2xq'e^{-x^{2}/2} - x^{2}qe^{-x^{2}/2} +q'e^{-x^{2}/2} + x^{2}qe^{-x^{2}/2}
	\end{align*}
	or
	\begin{align*}
		0 = -q'' +2xq' - (k+1)q
	\end{align*}
	
	Try a solution form $q(x) = \sum_{n=0}^{\infty} a_{n}x^{n}$ and derive a relationship for $a_{n+2}$ in terms of $a_{n}$ (Do not try to get a closed form for the coefficients). For what $k$ does the power series truncate (cut off) to be a polynomial? These are called the Hermite polynomials. 
	
	
	
	\textbf{Problem 3:} (Integrating the heat kernel) Recall that we used that the integral $\int_{\R^{n}} H_{t}(x) dx$ was finite during our proofs. In this exercise, we are going to verify this fact.
	
	Part A.) First, compute the Gaussian integral
	\begin{align*}
		I = \int_{-\infty}^{\infty} e^{-x^{2}} dx
	\end{align*} 
	\textit{Hint:} Try to instead compute $I^{2}$ by writing the integral as $I^2 \int_{-\infty}^{\infty}  \int_{-\infty}^{\infty} e^{-x^{2}-y^{2}} dy dx$. 
	
	Part B.) Use Part A to compute 
	\begin{align*}
		\int_{\R^{n}} e^{-|x|^{2}} dx
	\end{align*}
	
	Part C.) Use a change-of-variables to compute 
	\begin{align*}
		\int_{\R^{n}} (4\pi t)^{-n/2}e^{-|x|^{2}/4t} dx
	\end{align*}
	
	
	
	
	
	
	\textbf{Problem 4:} (Energy Estimates, Borthwick Problem 6.3 + 6.4) Recall that we used $u(t,x)$ to model the temperature when deriving the heat equation $\pd_{t}u - \Delta u = 0$. Assuming $u$ is a solution to this equation on $(0,\infty) \times \Omega$ for some bounded domain $\Omega \subset \R^{n}$, we define the total thermal energy 
	\begin{align*}
		U(t) = \int_{\Omega} u(t,x) dx
	\end{align*}
	
	Part A.) Assuming $u(t,x)$ satisfies Neumann boundary conditions 
	\begin{align*}
		\frac{\pd}{\pd \nu} |_{\pd \Omega} = 0
	\end{align*}
	(for normal vector $\nu$ to $\pd \Omega$), show $U(t)$ is constant. Interpret the meaning of this physically. 
	
	Part B.) Define the energy of the function
	\begin{align*}
		\eta(t) = \int_{\Omega} u(t,x)^{2} dx
	\end{align*}
	and assume $u(t,x) |_{\pd \Omega} = 0$. Show that $\eta$ decreases as a function of time. 
	
	Part C.) Use B.) to show that a solution $u(t,x)$ satisfying $u|_{t=0} = g$ and $u|_{\pd \Omega} = h$, for $g,h$ continuous on $\Omega$ and $\pd \Omega$, us uniquely determined by $g$ and $h$.
	
	
\end{document}
